% Tutorial set: 03
% Source: Tutorial Three, pp. 3-1--3-2 (PDF pp. 1--2), Questions 1--4
% Session date: 2026-09-01
% Human verified: Yes

\clearpage
\subsection{Tutorial 03：CPU Scheduling}
\label{tutorial:SC2005-TUT03}

\exercisemeta
  {SC2005-TUT03}
  {2026-09-01}
  {Tutorial Three，pp.~3-1--3-2（PDF pp.~1--2），Questions 1--4}
  {全部题目已作答并完成订正}
  {已确认（2026-09-01）}

\exercisequestion{SC2005-TUT03-Q01}{Scheduler 与 Process Migration 判断}

\textbf{对应讲义：}
\relatedtopic{SC2005-L05-T01}{CPU--I/O Bursts 与 Short-Term Scheduler}；
\relatedtopic{SC2005-L06-T05}{Partitioned 与 Global Multiprocessor Scheduling}

\subsubsection*{题目}

判断下列陈述是 True 还是 False，并说明理由。

\begin{enumerate}[label=(\alph*)]
  \item 当一个新 process（进程）被 admitted（接纳）进入系统时，为了保持 CPU busy，short-term scheduler（短期调度器）必须执行。
  \item Partitioned multiprocessor scheduling（分区式多处理器调度）会因 private core-specific cache（核心私有缓存）中的数据而承受 process migration overhead（进程迁移开销）。
\end{enumerate}

\begin{checkedsolution}
\begin{enumerate}[label=(\alph*)]
  \item \textbf{False。}新 process 进入 Ready state（就绪状态）不代表 CPU 当前 idle，也不代表必须立即改换 Running process。只有 CPU 需要选择下一个 process，或 preemptive policy（抢占式策略）在该事件上重新调度时，short-term scheduler 才需要介入。
  \item \textbf{False（按课程的严格 partitioned model）。}Process 创建后固定映射到某个 core，并进入该 core 的 local ready queue，因此普通执行中不会在 cores 之间迁移；其典型代价是 load imbalance（负载不均）。Migration 与 cache reload overhead 主要属于 global scheduling；若实现额外进行 load rebalancing，则另当别论。
\end{enumerate}
\end{checkedsolution}

\exercisequestion{SC2005-TUT03-Q02}{RR 与 Dual-core Global SRTF}

\textbf{对应讲义：}
\relatedtopic{SC2005-L05-T03}{Scheduling Objectives 与 Time Metrics}；
\relatedtopic{SC2005-L06-T01}{SJF 与 SRTF}；
\relatedtopic{SC2005-L06-T03}{Round Robin 与 Time Quantum}；
\relatedtopic{SC2005-L06-T05}{Partitioned 与 Global Multiprocessor Scheduling}

\subsubsection*{题目}

五个 processes 共享同一个 I/O device，参数如下。CPU burst 用逗号分隔；I/O 在第一个 CPU burst 后发生。

\begin{center}
\small
\begin{tabular}{@{}ccccc@{}}
  \toprule
  Process & CPU bursts & I/O burst & Arrival & 同时到达顺序 \\
  \midrule
  $P_1$ & $2,8$ & $2$ & $0$ & 1 \\
  $P_2$ & $1$   & $0$ & $0$ & 2 \\
  $P_3$ & $1,1$ & $1$ & $2$ & 1 \\
  $P_4$ & $1$   & $0$ & $2$ & 2 \\
  $P_5$ & $2,3$ & $2$ & $4$ & --- \\
  \bottomrule
\end{tabular}
\end{center}

\begin{enumerate}[label=(\alph*)]
  \item 分别画出 RR（Round Robin，时间片轮转）$q=2$ 的 single-core schedule，以及 two-core global SRTF schedule；
  \item 求每个 process 的 turnaround time（周转时间）和 waiting time（就绪等待时间）。
\end{enumerate}

\begin{checkedsolution}
图中采用一个明确的 simultaneous-event convention（同时事件约定）：RR 在 $t=4$ 先把新到达的 $P_5$ 加入 Ready queue，再加入刚完成 I/O 的 $P_1$。若课程采用另一种同时事件顺序，RR 的局部次序可能不同；下面的结果与本次课堂订正一致。

\begin{figure}[htbp]
  \centering
  \resizebox{0.98\textwidth}{!}{\begin{tikzpicture}[x=0.55cm,y=0.72cm,>=Latex,font=\scriptsize]
  \node[anchor=west,font=\small\bfseries] at (0,7.15) {RR, $q=2$, single core};

  \node[anchor=east] at (-0.3,6.25) {CPU};
  \draw[fill=blue!12]   (0,5.9) rectangle (2,6.6) node[midway] {$P_1$};
  \draw[fill=orange!18] (2,5.9) rectangle (3,6.6) node[midway] {$P_2$};
  \draw[fill=green!15]  (3,5.9) rectangle (4,6.6) node[midway] {$P_3$};
  \draw[fill=red!12]    (4,5.9) rectangle (5,6.6) node[midway] {$P_4$};
  \draw[fill=purple!15] (5,5.9) rectangle (7,6.6) node[midway] {$P_5$};
  \draw[fill=blue!12]   (7,5.9) rectangle (9,6.6) node[midway] {$P_1$};
  \draw[fill=green!15]  (9,5.9) rectangle (10,6.6) node[midway] {$P_3$};
  \draw[fill=blue!12]   (10,5.9) rectangle (12,6.6) node[midway] {$P_1$};
  \draw[fill=purple!15] (12,5.9) rectangle (14,6.6) node[midway] {$P_5$};
  \draw[fill=blue!12]   (14,5.9) rectangle (16,6.6) node[midway] {$P_1$};
  \draw[fill=purple!15] (16,5.9) rectangle (17,6.6) node[midway] {$P_5$};
  \draw[fill=blue!12]   (17,5.9) rectangle (19,6.6) node[midway] {$P_1$};

  \node[anchor=east] at (-0.3,5.15) {I/O};
  \draw[fill=blue!12]   (2,4.8) rectangle (4,5.5) node[midway] {$P_1$};
  \draw[fill=green!15]  (4,4.8) rectangle (5,5.5) node[midway] {$P_3$};
  \draw[fill=purple!15] (7,4.8) rectangle (9,5.5) node[midway] {$P_5$};
  \foreach \t in {0,2,3,4,5,7,9,10,12,14,16,17,19}{
    \draw (\t,4.72)--(\t,4.58) node[below,font=\tiny]{\t};
  }

  \node[anchor=west,font=\small\bfseries] at (0,3.65) {Global SRTF, two cores};

  \node[anchor=east] at (-0.3,2.75) {Core 1};
  \draw[fill=blue!12]  (0,2.4) rectangle (2,3.1) node[midway] {$P_1$};
  \draw[fill=red!12]   (2,2.4) rectangle (3,3.1) node[midway] {$P_4$};
  \draw[dashed]        (3,2.4) rectangle (4,3.1) node[midway,font=\tiny] {idle};
  \draw[fill=blue!12]  (4,2.4) rectangle (5,3.1) node[midway] {$P_1$};
  \draw[fill=green!15] (5,2.4) rectangle (6,3.1) node[midway] {$P_3$};
  \draw[fill=blue!12]  (6,2.4) rectangle (13,3.1) node[midway] {$P_1$};

  \node[anchor=east] at (-0.3,1.75) {Core 2};
  \draw[fill=orange!18] (0,1.4) rectangle (1,2.1) node[midway] {$P_2$};
  \draw[dashed]         (1,1.4) rectangle (2,2.1) node[midway,font=\tiny] {idle};
  \draw[fill=green!15]  (2,1.4) rectangle (3,2.1) node[midway] {$P_3$};
  \draw[dashed]         (3,1.4) rectangle (4,2.1) node[midway,font=\tiny] {idle};
  \draw[fill=purple!15] (4,1.4) rectangle (6,2.1) node[midway] {$P_5$};
  \draw[dashed]         (6,1.4) rectangle (8,2.1) node[midway,font=\tiny] {idle};
  \draw[fill=purple!15] (8,1.4) rectangle (11,2.1) node[midway] {$P_5$};

  \node[anchor=east] at (-0.3,0.75) {I/O};
  \draw[fill=blue!12]   (2,0.4) rectangle (4,1.1) node[midway] {$P_1$};
  \draw[fill=green!15]  (4,0.4) rectangle (5,1.1) node[midway] {$P_3$};
  \draw[fill=purple!15] (6,0.4) rectangle (8,1.1) node[midway] {$P_5$};
  \foreach \t in {0,1,2,3,4,5,6,8,11,13}{
    \draw (\t,0.32)--(\t,0.18) node[below,font=\tiny]{\t};
  }
\end{tikzpicture}
}
  \caption{Question 2 的 RR 与 global SRTF Gantt charts；I/O 行只画实际占用共享设备的区间。}
  \label{fig:sc2005-tut03-q2-schedules}
\end{figure}

令 $C_i$ 为 completion time，则
\[
  \text{turnaround}_i=C_i-A_i,
  \qquad
  \text{waiting}_i=\text{Ready queue 中累计等待时间}.
\]

\begin{center}
\small
\begin{tabular}{@{}c|ccc|ccc@{}}
  \toprule
  & \multicolumn{3}{c|}{RR, $q=2$} & \multicolumn{3}{c}{Global SRTF, two cores} \\
  Process & $C_i$ & Turnaround & Waiting & $C_i$ & Turnaround & Waiting \\
  \midrule
  $P_1$ & 19 & 19 & 7 & 13 & 13 & 1 \\
  $P_2$ & 3  & 3  & 2 & 1  & 1  & 0 \\
  $P_3$ & 10 & 8  & 5 & 6  & 4  & 0 \\
  $P_4$ & 5  & 3  & 2 & 3  & 1  & 0 \\
  $P_5$ & 17 & 13 & 6 & 11 & 7  & 0 \\
  \bottomrule
\end{tabular}
\end{center}

SRTF 中 $P_3$ 在 $t=3$ 完成第一个 CPU burst，但共享 I/O device 被 $P_1$ 占用到 $t=4$，所以它在 $t=3$--$4$ 等待设备、$t=4$--$5$ 执行 I/O。该一单位时间不是 Ready-queue waiting，因此 $P_3$ 的 CPU waiting time 仍为 $0$；此时不能直接用 ``turnaround $-$ CPU bursts $-$ I/O service time'' 计算 waiting，否则会把 I/O-device queueing 混入结果。
\end{checkedsolution}

\subsubsection*{检查与常见错误}

SRTF 比较的是当前 CPU burst 的 remaining time，而不是 process 余下所有 CPU 与 I/O 时间之和。多个 cores 可以同时运行不同 processes，但同一个 process 在任一时刻至多占用一个 core。

\exercisequestion{SC2005-TUT03-Q03}{由 Dual-core Schedule 反推 Bursts 与 Policy}

\textbf{对应讲义：}
\begin{itemize}[leftmargin=2em]
  \item \relatedtopic{SC2005-L06-T01}{SJF 与 SRTF}
  \item \relatedtopic{SC2005-L06-T05}{Partitioned 与 Global Multiprocessor Scheduling}
\end{itemize}

\subsubsection*{题目}

$P_0,P_1,P_4,P_5,P_6$ 在 $t=0$ 依次到达，$P_2$ 在 $t=3$ 到达，$P_3$ 在 $t=9$ 到达。根据给定 dual-core schedule：

\begin{enumerate}[label=(\alph*)]
  \item 反推出各 process 的全部 CPU burst 与 I/O burst lengths；
  \item 判断它属于 partitioned 还是 global multiprocessor scheduling，并从课堂算法中找出与图示相容的 algorithm，说明依据。
\end{enumerate}

\begin{figure}[htbp]
  \centering
  \resizebox{0.92\textwidth}{!}{\begin{tikzpicture}[x=0.74cm,y=0.82cm,>=Latex,font=\small]
  \node[anchor=east] at (-0.25,2.75) {I/O};
  \draw[fill=blue!12]   (1,2.4) rectangle (3,3.1) node[midway] {$P_0$};
  \draw[fill=red!12]    (4,2.4) rectangle (7,3.1) node[midway] {$P_4$};
  \draw[fill=orange!18] (7,2.4) rectangle (9,3.1) node[midway] {$P_2$};

  \node[anchor=east] at (-0.25,1.75) {Core 1};
  \draw[fill=blue!12]   (0,1.4) rectangle (1,2.1) node[midway] {$P_0$};
  \draw[fill=red!12]    (1,1.4) rectangle (4,2.1) node[midway] {$P_4$};
  \draw[fill=orange!18] (4,1.4) rectangle (5,2.1) node[midway] {$P_2$};
  \draw[fill=blue!12]   (5,1.4) rectangle (7,2.1) node[midway] {$P_0$};
  \draw[fill=red!12]    (7,1.4) rectangle (9,2.1) node[midway] {$P_4$};
  \draw[fill=orange!18] (9,1.4) rectangle (10,2.1) node[midway] {$P_2$};
  \draw[fill=green!15]  (10,1.4) rectangle (12,2.1) node[midway] {$P_3$};

  \node[anchor=east] at (-0.25,0.75) {Core 2};
  \draw[fill=purple!15] (0,0.4) rectangle (2,1.1) node[midway] {$P_5$};
  \draw[fill=teal!15]   (2,0.4) rectangle (6,1.1) node[midway] {$P_6$};
  \draw[fill=gray!18]   (6,0.4) rectangle (12,1.1) node[midway] {$P_1$};

  \foreach \t in {0,1,2,3,4,5,6,7,9,10,12}{
    \draw[red!70!black] (\t,0.32)--(\t,0.16) node[below,font=\scriptsize,text=black]{\t};
  }
\end{tikzpicture}
}
  \caption{Tutorial Three Question 3 所给 schedule 的等价重绘；时间区间与原图一致。}
  \label{fig:sc2005-tut03-q3-schedule}
\end{figure}

\begin{checkedsolution}
同一 CPU burst 可能被 preemption（抢占）切成多个 execution intervals；本图没有这种已知的切分，因此按连续执行区间读取：

\begin{center}
\small
\begin{tabular}{@{}ccl@{}}
  \toprule
  Process & CPU burst lengths & I/O burst lengths \\
  \midrule
  $P_0$ & $1,2$ & $2$（$t=1$--$3$） \\
  $P_1$ & $6$   & --- \\
  $P_2$ & $1,1$ & $2$（$t=7$--$9$；$t=5$--$7$ 等设备） \\
  $P_3$ & $2$   & --- \\
  $P_4$ & $3,2$ & $3$（$t=4$--$7$） \\
  $P_5$ & $2$   & --- \\
  $P_6$ & $4$   & --- \\
  \bottomrule
\end{tabular}
\end{center}

仅凭该 trace（轨迹）\textbf{不能唯一判断} partitioned 或 global：

\begin{itemize}
  \item 它与 \textbf{partitioned scheduling + per-core SJF} 相容；由于没有“严格更短的新 burst 到达并迫使抢占”的区间，per-core SRTF 也可能产生同一 trace；
  \item 它也与 \textbf{global nonpreemptive SJF} 相容：每次有 core 空闲时，从全局 Ready processes 中选择最短的 next CPU burst；
  \item 它与 \textbf{global SRTF 不相容}：$t=3$ 时新到达的 $P_2$ remaining time 为 $1$，而正在 Core 2 运行的 $P_6$ remaining time 为 $3$。Global SRTF 应让 $P_2$ 取代 $P_6$ 占用一个 core，但图中 $P_6$ 连续运行到 $t=6$。
\end{itemize}

因此，若题目要求唯一预期答案，最自然的课程解释是 partitioned scheduling，两个 local queues 使用 SJF；严格来说，图中信息不足以排除上述其他相容解释。
\end{checkedsolution}

\clearpage
\exercisequestion{SC2005-TUT03-Q04}{RR Quantum 与 CPU Efficiency}

\textbf{对应讲义：}
\relatedtopic{SC2005-L06-T03}{Round Robin 与 Time Quantum}

\subsubsection*{题目}

一个 process 在 blocking I/O 前平均运行 $T$，每次 process switch（进程切换）耗时 $S$，且该时间全部属于 overhead。定义 CPU efficiency，并分别求 RR 在以下条件下的效率：

\[
  \text{(a) }Q=\infty,\qquad
  \text{(b) }Q>T,\qquad
  \text{(c) }S<Q<T,\qquad
  \text{(d) }Q=S,\qquad
  \text{(e) }Q=0.
\]

\begin{checkedsolution}
CPU efficiency（CPU 效率）是总时间中用于执行 process instructions（有效工作）的比例：
\[
  \eta=\frac{\text{useful CPU execution time}}
             {\text{useful execution time}+\text{switch overhead time}}.
\]

\begin{center}
\small
\begin{tabularx}{0.94\textwidth}{@{}cclX@{}}
  \toprule
  Case & 条件 & Efficiency & 原因 \\
  \midrule
  (a) & $Q=\infty$ & $\displaystyle\frac{T}{T+S}$ & Process 运行 $T$ 后因 I/O blocking，再付出一次 switch $S$ \\
  (b) & $Q>T$ & $\displaystyle\frac{T}{T+S}$ & Blocking 先于 quantum expiry，行为与 (a) 相同 \\
  (c) & $S<Q<T$ & $\displaystyle\frac{Q}{Q+S}$ & 每段有效执行约为 $Q$，随后付出 switch $S$ \\
  (d) & $Q=S$ & $\displaystyle\frac{1}{2}$ & 有效执行与 overhead 等长 \\
  (e) & $Q=0$ & $0$ & 视为 $Q\to0^+$ 的极限；没有时间完成有效工作 \\
  \bottomrule
\end{tabularx}
\end{center}

Case (c) 是 steady-state（稳态）表达式；若 $T$ 不是 $Q$ 的整数倍，最后一段会短于 $Q$，精确的单次-burst 比例还取决于 switch 的计数约定。考试中的标准答案采用 $Q/(Q+S)$。
\end{checkedsolution}

\subsubsection*{检查与常见错误}

分母必须包含 useful time 与 overhead，故 (a)、(b) 写作 $T/(T+S)$，不能省略括号写成 $T/T+S$。$Q=0$ 不是可实际运行的配置，而是用于观察 quantum 趋零时 overhead 吞噬 CPU 的极限情形。
